% New slide.

\begin{frame}[t, fragile]{Haskell function which performs IO and returns a value}

\begin{minted}[linenos]{haskell}
addTwoNumbers_printAndReturnResult :: Int -> Int -> IO Int
addTwoNumbers_printAndReturnResult arg1 arg2 = do

    let

        sum = arg1 + arg2


    printf "Sum = %d" sum

    return sum
\end{minted}

\end{frame}


% New slide.

\begin{frame}[t, fragile]{Using parametric polymorphism to define a function}

Consider the function addTwoNumbers from earlier. Recall that it could only receive values of type Int and that it
would only ever return a value of type Int. But what if instead of just operating on values of type Int, we wanted it to
operate on values of any numerical type?

In that case, we could redefine our function to something like the following;

\begin{minted}[linenos]{haskell}
addTwoNumbers :: Num a => a -> a -> a
addTwoNumbers    arg1 arg2 =

    let

        sum = arg1 + arg2

    in

        sum
\end{minted}

The key here is the revised declaration for the function, that is;

\begin{minted}[linenos=false]{haskell}
addTwoNumbers :: Num a => a -> a -> a
\end{minted}

This new declaration informs us that the function can now receive arguments of any type `a', so long as these types
derive from the type Num.

\end{frame}


% New slide.

\begin{frame}[t, fragile]{Type variables}

Alternatively, we could have declared the function addTwoNumbers as follows;

\begin{minted}[linenos=false]{haskell}
addTwoNumbers :: a -> a -> a
\end{minted}

where `a' denotes a type variable, which stands for any type of variable.

But what if `a' stood for variables of type Char, so that we essentially had;

\begin{minted}[linenos=false]{haskell}
addTwoNumbers :: Char a => a -> a -> a
\end{minted}

We could maybe think of a way to make this make sense. But what if we had the following function declaration;

\begin{minted}[linenos=false]{haskell}
convertToUpperCase :: a -> a
\end{minted}

We can infer from this function declaration -- or, to borrow from the  Haskell parlance, function type signature, that
it is meant to take arguments of type Char. So what then would happen if we tried to invoke it as follows?

\begin{minted}[linenos=false]{haskell}
convertToUpperCase (3.141 :: Float)
\end{minted}

Nothing good would probably come of it.

\end{frame}


% New slide.

\begin{frame}[t, fragile]{Type variables}

To prevent problems like this from happening in the first place, we could amend the function so that it only takes
arguments which are of type Char.

\begin{minted}[linenos=false]{haskell}
convertToUpperCase :: Char a => a -> a -> a
\end{minted}


\end{frame}